% ======================================================================
% TỔNG HỢP KIẾN THỨC CHƯƠNG 1--4 - LÝ THUYẾT ĐIỀU KHIỂN TỰ ĐỘNG
% LaTeX fragment để chèn sau vào main.tex bằng:
%   % ======================================================================
% TỔNG HỢP KIẾN THỨC CHƯƠNG 1--4 - LÝ THUYẾT ĐIỀU KHIỂN TỰ ĐỘNG
% LaTeX fragment để chèn sau vào main.tex bằng:
%   % ======================================================================
% TỔNG HỢP KIẾN THỨC CHƯƠNG 1--4 - LÝ THUYẾT ĐIỀU KHIỂN TỰ ĐỘNG
% LaTeX fragment để chèn sau vào main.tex bằng:
%   % ======================================================================
% TỔNG HỢP KIẾN THỨC CHƯƠNG 1--4 - LÝ THUYẾT ĐIỀU KHIỂN TỰ ĐỘNG
% LaTeX fragment để chèn sau vào main.tex bằng:
%   \input{tong_hop_kien_thuc_chuong_1_4}
% Không có \documentclass, \begin{document}, \end{document}.
% ======================================================================

\providecommand{\Lap}{\mathcal{L}}
\providecommand{\dd}{\,\mathrm{d}}
\providecommand{\e}{\mathrm{e}}

\clearpage
\section*{TỔNG HỢP KIẾN THỨC VÀ PHƯƠNG PHÁP GIẢI BÀI TẬP -- CHƯƠNG 1 ĐẾN CHƯƠNG 4}

\noindent
Phần này được viết theo hướng tra nhanh khi làm bài: mỗi chương gồm kiến thức cốt lõi,
công thức cần nhớ, quy trình giải các dạng bài thường gặp, ví dụ mẫu và lỗi hay mắc.
Ký hiệu dùng thống nhất:
\[
r(t):\text{ tín hiệu đặt},\qquad
e(t):\text{ sai lệch},\qquad
u(t):\text{ tín hiệu điều khiển},\qquad
c(t)\text{ hoặc }y(t):\text{ tín hiệu ra}.
\]

% ======================================================================
\subsection*{CHƯƠNG 1 -- TỔNG QUAN VỀ HỆ THỐNG ĐIỀU KHIỂN TỰ ĐỘNG}
% ======================================================================

\subsubsection*{1. Khái niệm cốt lõi}

Điều khiển là quá trình \textbf{thu thập thông tin, xử lý thông tin và tác động lên hệ thống}
để đáp ứng thực tế gần với mục đích định trước. Điều khiển tự động là quá trình điều khiển
không cần con người trực tiếp can thiệp trong quá trình vận hành.

\begin{center}
\begin{tabular}{|p{0.20\linewidth}|p{0.27\linewidth}|p{0.43\linewidth}|}
\hline
\textbf{Ký hiệu} & \textbf{Tên} & \textbf{Ý nghĩa} \\ \hline
$r(t)$ & Tín hiệu chuẩn / tín hiệu đặt & Giá trị mong muốn của hệ thống \\ \hline
$c(t)$, $y(t)$ & Tín hiệu ra & Giá trị thực tế của đại lượng cần điều khiển \\ \hline
$c_{ht}(t)$ & Tín hiệu hồi tiếp & Tín hiệu đo được đưa trở lại bộ so sánh \\ \hline
$e(t)$ & Sai lệch & Thường có dạng $e(t)=r(t)-c_{ht}(t)$ với hồi tiếp âm \\ \hline
$u(t)$ & Tín hiệu điều khiển & Tín hiệu từ bộ điều khiển tác động lên đối tượng \\ \hline
\end{tabular}
\end{center}

\subsubsection*{2. Hệ vòng hở và hệ vòng kín}

\textbf{Hệ vòng hở (open-loop):} không dùng tín hiệu ra để hiệu chỉnh tác động điều khiển:
\[
r(t)\longrightarrow \text{Bộ điều khiển}\longrightarrow \text{Đối tượng}\longrightarrow c(t).
\]

\textbf{Hệ vòng kín (closed-loop):} dùng tín hiệu phản hồi để so sánh với tín hiệu đặt,
tạo sai lệch và hiệu chỉnh liên tục:
\[
e(t)=r(t)-c_{ht}(t).
\]

Gợi ý nhận biết: nếu sơ đồ có đường từ đầu ra quay trở lại nút so sánh thì đó là
hệ vòng kín; nếu không có đường phản hồi thì là hệ vòng hở.

\subsubsection*{3. Các bài toán điều khiển cơ bản}

\begin{itemize}
  \item \textbf{Phân tích hệ thống:} đã biết cấu trúc và thông số; cần tìm đáp ứng, đặc tính và đánh giá chất lượng.
  \item \textbf{Thiết kế hệ thống:} đã biết đối tượng; cần thiết kế bộ điều khiển để đạt yêu cầu.
  \item \textbf{Nhận dạng hệ thống:} chưa biết đầy đủ cấu trúc/thông số; dùng dữ liệu vào--ra để xác định mô hình.
\end{itemize}

\subsubsection*{4. Các nguyên tắc và phân loại thường gặp}

Trong tài liệu, các nguyên tắc được nhấn mạnh gồm thông tin phản hồi, đa dạng tương xứng,
bổ sung ngoài, dự trữ, phân cấp và cân bằng nội. Ba cấu trúc điều khiển theo thông tin phản hồi
thường gặp là bù nhiễu, san bằng sai lệch và phối hợp hai nguyên tắc trên.

Các cách phân loại chính:
\begin{itemize}
  \item Tuyến tính / phi tuyến.
  \item Bất biến theo thời gian / biến đổi theo thời gian.
  \item Liên tục / rời rạc / số.
  \item SISO (một vào--một ra) / MIMO (nhiều vào--nhiều ra).
  \item Theo mục tiêu: ổn định hoá, theo chương trình, bám theo tín hiệu, thích nghi, tối ưu.
\end{itemize}

\subsubsection*{5. Quy trình thiết lập một hệ thống điều khiển}

Có thể ghi nhớ theo chuỗi:
\[
\boxed{
\text{Yêu cầu kỹ thuật}
\rightarrow
\text{Mô tả hệ vật lý}
\rightarrow
\text{Sơ đồ chức năng}
\rightarrow
\text{Mô hình toán}
\rightarrow
\text{Sơ đồ khối/graph/KGTT}
\rightarrow
\text{Rút gọn}
\rightarrow
\text{Phân tích, thiết kế}
}
\]

\subsubsection*{6. Dạng bài Chương 1: nhận dạng thành phần của hệ}

\textbf{Quy trình:}
\begin{enumerate}
  \item Xác định đại lượng mong muốn $\Rightarrow r(t)$.
  \item Xác định đại lượng thực tế cần điều khiển $\Rightarrow c(t)$.
  \item Xác định thiết bị đo tạo tín hiệu phản hồi.
  \item Xác định bộ điều khiển.
  \item Xác định cơ cấu chấp hành.
  \item Xác định đối tượng điều khiển.
  \item Kiểm tra có đường phản hồi hay không để kết luận vòng hở/vòng kín.
\end{enumerate}

\textbf{Ví dụ mẫu -- điều khiển mức chất lỏng.}
Mức mong muốn là $r(t)$; cảm biến/phao đo mức tạo phản hồi; bộ điều khiển tính sai lệch;
van là cơ cấu chấp hành; bồn là đối tượng; mức thực tế là $c(t)$. Nếu mức đo được quay về
bộ so sánh thì hệ là vòng kín.

\medskip
\noindent
\fbox{\parbox{0.96\linewidth}{
\textbf{Lỗi hay gặp Chương 1:}
nhầm tín hiệu điều khiển $u(t)$ với tín hiệu đặt $r(t)$;
nhầm cảm biến với bộ điều khiển; thấy có cảm biến nhưng không kiểm tra xem tín hiệu đo
có thực sự quay về bộ so sánh hay không.
}}


% Không có \documentclass, \begin{document}, \end{document}.
% ======================================================================

\providecommand{\Lap}{\mathcal{L}}
\providecommand{\dd}{\,\mathrm{d}}
\providecommand{\e}{\mathrm{e}}

\clearpage
\section*{TỔNG HỢP KIẾN THỨC VÀ PHƯƠNG PHÁP GIẢI BÀI TẬP -- CHƯƠNG 1 ĐẾN CHƯƠNG 4}

\noindent
Phần này được viết theo hướng tra nhanh khi làm bài: mỗi chương gồm kiến thức cốt lõi,
công thức cần nhớ, quy trình giải các dạng bài thường gặp, ví dụ mẫu và lỗi hay mắc.
Ký hiệu dùng thống nhất:
\[
r(t):\text{ tín hiệu đặt},\qquad
e(t):\text{ sai lệch},\qquad
u(t):\text{ tín hiệu điều khiển},\qquad
c(t)\text{ hoặc }y(t):\text{ tín hiệu ra}.
\]

% ======================================================================
\subsection*{CHƯƠNG 1 -- TỔNG QUAN VỀ HỆ THỐNG ĐIỀU KHIỂN TỰ ĐỘNG}
% ======================================================================

\subsubsection*{1. Khái niệm cốt lõi}

Điều khiển là quá trình \textbf{thu thập thông tin, xử lý thông tin và tác động lên hệ thống}
để đáp ứng thực tế gần với mục đích định trước. Điều khiển tự động là quá trình điều khiển
không cần con người trực tiếp can thiệp trong quá trình vận hành.

\begin{center}
\begin{tabular}{|p{0.20\linewidth}|p{0.27\linewidth}|p{0.43\linewidth}|}
\hline
\textbf{Ký hiệu} & \textbf{Tên} & \textbf{Ý nghĩa} \\ \hline
$r(t)$ & Tín hiệu chuẩn / tín hiệu đặt & Giá trị mong muốn của hệ thống \\ \hline
$c(t)$, $y(t)$ & Tín hiệu ra & Giá trị thực tế của đại lượng cần điều khiển \\ \hline
$c_{ht}(t)$ & Tín hiệu hồi tiếp & Tín hiệu đo được đưa trở lại bộ so sánh \\ \hline
$e(t)$ & Sai lệch & Thường có dạng $e(t)=r(t)-c_{ht}(t)$ với hồi tiếp âm \\ \hline
$u(t)$ & Tín hiệu điều khiển & Tín hiệu từ bộ điều khiển tác động lên đối tượng \\ \hline
\end{tabular}
\end{center}

\subsubsection*{2. Hệ vòng hở và hệ vòng kín}

\textbf{Hệ vòng hở (open-loop):} không dùng tín hiệu ra để hiệu chỉnh tác động điều khiển:
\[
r(t)\longrightarrow \text{Bộ điều khiển}\longrightarrow \text{Đối tượng}\longrightarrow c(t).
\]

\textbf{Hệ vòng kín (closed-loop):} dùng tín hiệu phản hồi để so sánh với tín hiệu đặt,
tạo sai lệch và hiệu chỉnh liên tục:
\[
e(t)=r(t)-c_{ht}(t).
\]

Gợi ý nhận biết: nếu sơ đồ có đường từ đầu ra quay trở lại nút so sánh thì đó là
hệ vòng kín; nếu không có đường phản hồi thì là hệ vòng hở.

\subsubsection*{3. Các bài toán điều khiển cơ bản}

\begin{itemize}
  \item \textbf{Phân tích hệ thống:} đã biết cấu trúc và thông số; cần tìm đáp ứng, đặc tính và đánh giá chất lượng.
  \item \textbf{Thiết kế hệ thống:} đã biết đối tượng; cần thiết kế bộ điều khiển để đạt yêu cầu.
  \item \textbf{Nhận dạng hệ thống:} chưa biết đầy đủ cấu trúc/thông số; dùng dữ liệu vào--ra để xác định mô hình.
\end{itemize}

\subsubsection*{4. Các nguyên tắc và phân loại thường gặp}

Trong tài liệu, các nguyên tắc được nhấn mạnh gồm thông tin phản hồi, đa dạng tương xứng,
bổ sung ngoài, dự trữ, phân cấp và cân bằng nội. Ba cấu trúc điều khiển theo thông tin phản hồi
thường gặp là bù nhiễu, san bằng sai lệch và phối hợp hai nguyên tắc trên.

Các cách phân loại chính:
\begin{itemize}
  \item Tuyến tính / phi tuyến.
  \item Bất biến theo thời gian / biến đổi theo thời gian.
  \item Liên tục / rời rạc / số.
  \item SISO (một vào--một ra) / MIMO (nhiều vào--nhiều ra).
  \item Theo mục tiêu: ổn định hoá, theo chương trình, bám theo tín hiệu, thích nghi, tối ưu.
\end{itemize}

\subsubsection*{5. Quy trình thiết lập một hệ thống điều khiển}

Có thể ghi nhớ theo chuỗi:
\[
\boxed{
\text{Yêu cầu kỹ thuật}
\rightarrow
\text{Mô tả hệ vật lý}
\rightarrow
\text{Sơ đồ chức năng}
\rightarrow
\text{Mô hình toán}
\rightarrow
\text{Sơ đồ khối/graph/KGTT}
\rightarrow
\text{Rút gọn}
\rightarrow
\text{Phân tích, thiết kế}
}
\]

\subsubsection*{6. Dạng bài Chương 1: nhận dạng thành phần của hệ}

\textbf{Quy trình:}
\begin{enumerate}
  \item Xác định đại lượng mong muốn $\Rightarrow r(t)$.
  \item Xác định đại lượng thực tế cần điều khiển $\Rightarrow c(t)$.
  \item Xác định thiết bị đo tạo tín hiệu phản hồi.
  \item Xác định bộ điều khiển.
  \item Xác định cơ cấu chấp hành.
  \item Xác định đối tượng điều khiển.
  \item Kiểm tra có đường phản hồi hay không để kết luận vòng hở/vòng kín.
\end{enumerate}

\textbf{Ví dụ mẫu -- điều khiển mức chất lỏng.}
Mức mong muốn là $r(t)$; cảm biến/phao đo mức tạo phản hồi; bộ điều khiển tính sai lệch;
van là cơ cấu chấp hành; bồn là đối tượng; mức thực tế là $c(t)$. Nếu mức đo được quay về
bộ so sánh thì hệ là vòng kín.

\medskip
\noindent
\fbox{\parbox{0.96\linewidth}{
\textbf{Lỗi hay gặp Chương 1:}
nhầm tín hiệu điều khiển $u(t)$ với tín hiệu đặt $r(t)$;
nhầm cảm biến với bộ điều khiển; thấy có cảm biến nhưng không kiểm tra xem tín hiệu đo
có thực sự quay về bộ so sánh hay không.
}}


% Không có \documentclass, \begin{document}, \end{document}.
% ======================================================================

\providecommand{\Lap}{\mathcal{L}}
\providecommand{\dd}{\,\mathrm{d}}
\providecommand{\e}{\mathrm{e}}

\clearpage
\section*{TỔNG HỢP KIẾN THỨC VÀ PHƯƠNG PHÁP GIẢI BÀI TẬP -- CHƯƠNG 1 ĐẾN CHƯƠNG 4}

\noindent
Phần này được viết theo hướng tra nhanh khi làm bài: mỗi chương gồm kiến thức cốt lõi,
công thức cần nhớ, quy trình giải các dạng bài thường gặp, ví dụ mẫu và lỗi hay mắc.
Ký hiệu dùng thống nhất:
\[
r(t):\text{ tín hiệu đặt},\qquad
e(t):\text{ sai lệch},\qquad
u(t):\text{ tín hiệu điều khiển},\qquad
c(t)\text{ hoặc }y(t):\text{ tín hiệu ra}.
\]

% ======================================================================
\subsection*{CHƯƠNG 1 -- TỔNG QUAN VỀ HỆ THỐNG ĐIỀU KHIỂN TỰ ĐỘNG}
% ======================================================================

\subsubsection*{1. Khái niệm cốt lõi}

Điều khiển là quá trình \textbf{thu thập thông tin, xử lý thông tin và tác động lên hệ thống}
để đáp ứng thực tế gần với mục đích định trước. Điều khiển tự động là quá trình điều khiển
không cần con người trực tiếp can thiệp trong quá trình vận hành.

\begin{center}
\begin{tabular}{|p{0.20\linewidth}|p{0.27\linewidth}|p{0.43\linewidth}|}
\hline
\textbf{Ký hiệu} & \textbf{Tên} & \textbf{Ý nghĩa} \\ \hline
$r(t)$ & Tín hiệu chuẩn / tín hiệu đặt & Giá trị mong muốn của hệ thống \\ \hline
$c(t)$, $y(t)$ & Tín hiệu ra & Giá trị thực tế của đại lượng cần điều khiển \\ \hline
$c_{ht}(t)$ & Tín hiệu hồi tiếp & Tín hiệu đo được đưa trở lại bộ so sánh \\ \hline
$e(t)$ & Sai lệch & Thường có dạng $e(t)=r(t)-c_{ht}(t)$ với hồi tiếp âm \\ \hline
$u(t)$ & Tín hiệu điều khiển & Tín hiệu từ bộ điều khiển tác động lên đối tượng \\ \hline
\end{tabular}
\end{center}

\subsubsection*{2. Hệ vòng hở và hệ vòng kín}

\textbf{Hệ vòng hở (open-loop):} không dùng tín hiệu ra để hiệu chỉnh tác động điều khiển:
\[
r(t)\longrightarrow \text{Bộ điều khiển}\longrightarrow \text{Đối tượng}\longrightarrow c(t).
\]

\textbf{Hệ vòng kín (closed-loop):} dùng tín hiệu phản hồi để so sánh với tín hiệu đặt,
tạo sai lệch và hiệu chỉnh liên tục:
\[
e(t)=r(t)-c_{ht}(t).
\]

Gợi ý nhận biết: nếu sơ đồ có đường từ đầu ra quay trở lại nút so sánh thì đó là
hệ vòng kín; nếu không có đường phản hồi thì là hệ vòng hở.

\subsubsection*{3. Các bài toán điều khiển cơ bản}

\begin{itemize}
  \item \textbf{Phân tích hệ thống:} đã biết cấu trúc và thông số; cần tìm đáp ứng, đặc tính và đánh giá chất lượng.
  \item \textbf{Thiết kế hệ thống:} đã biết đối tượng; cần thiết kế bộ điều khiển để đạt yêu cầu.
  \item \textbf{Nhận dạng hệ thống:} chưa biết đầy đủ cấu trúc/thông số; dùng dữ liệu vào--ra để xác định mô hình.
\end{itemize}

\subsubsection*{4. Các nguyên tắc và phân loại thường gặp}

Trong tài liệu, các nguyên tắc được nhấn mạnh gồm thông tin phản hồi, đa dạng tương xứng,
bổ sung ngoài, dự trữ, phân cấp và cân bằng nội. Ba cấu trúc điều khiển theo thông tin phản hồi
thường gặp là bù nhiễu, san bằng sai lệch và phối hợp hai nguyên tắc trên.

Các cách phân loại chính:
\begin{itemize}
  \item Tuyến tính / phi tuyến.
  \item Bất biến theo thời gian / biến đổi theo thời gian.
  \item Liên tục / rời rạc / số.
  \item SISO (một vào--một ra) / MIMO (nhiều vào--nhiều ra).
  \item Theo mục tiêu: ổn định hoá, theo chương trình, bám theo tín hiệu, thích nghi, tối ưu.
\end{itemize}

\subsubsection*{5. Quy trình thiết lập một hệ thống điều khiển}

Có thể ghi nhớ theo chuỗi:
\[
\boxed{
\text{Yêu cầu kỹ thuật}
\rightarrow
\text{Mô tả hệ vật lý}
\rightarrow
\text{Sơ đồ chức năng}
\rightarrow
\text{Mô hình toán}
\rightarrow
\text{Sơ đồ khối/graph/KGTT}
\rightarrow
\text{Rút gọn}
\rightarrow
\text{Phân tích, thiết kế}
}
\]

\subsubsection*{6. Dạng bài Chương 1: nhận dạng thành phần của hệ}

\textbf{Quy trình:}
\begin{enumerate}
  \item Xác định đại lượng mong muốn $\Rightarrow r(t)$.
  \item Xác định đại lượng thực tế cần điều khiển $\Rightarrow c(t)$.
  \item Xác định thiết bị đo tạo tín hiệu phản hồi.
  \item Xác định bộ điều khiển.
  \item Xác định cơ cấu chấp hành.
  \item Xác định đối tượng điều khiển.
  \item Kiểm tra có đường phản hồi hay không để kết luận vòng hở/vòng kín.
\end{enumerate}

\textbf{Ví dụ mẫu -- điều khiển mức chất lỏng.}
Mức mong muốn là $r(t)$; cảm biến/phao đo mức tạo phản hồi; bộ điều khiển tính sai lệch;
van là cơ cấu chấp hành; bồn là đối tượng; mức thực tế là $c(t)$. Nếu mức đo được quay về
bộ so sánh thì hệ là vòng kín.

\medskip
\noindent
\fbox{\parbox{0.96\linewidth}{
\textbf{Lỗi hay gặp Chương 1:}
nhầm tín hiệu điều khiển $u(t)$ với tín hiệu đặt $r(t)$;
nhầm cảm biến với bộ điều khiển; thấy có cảm biến nhưng không kiểm tra xem tín hiệu đo
có thực sự quay về bộ so sánh hay không.
}}


% Không có \documentclass, \begin{document}, \end{document}.
% ======================================================================

\providecommand{\Lap}{\mathcal{L}}
\providecommand{\dd}{\,\mathrm{d}}
\providecommand{\e}{\mathrm{e}}

\clearpage
\section*{TỔNG HỢP KIẾN THỨC VÀ PHƯƠNG PHÁP GIẢI BÀI TẬP -- CHƯƠNG 1 ĐẾN CHƯƠNG 4}

\noindent
Phần này được viết theo hướng tra nhanh khi làm bài: mỗi chương gồm kiến thức cốt lõi,
công thức cần nhớ, quy trình giải các dạng bài thường gặp, ví dụ mẫu và lỗi hay mắc.
Ký hiệu dùng thống nhất:
\[
r(t):\text{ tín hiệu đặt},\qquad
e(t):\text{ sai lệch},\qquad
u(t):\text{ tín hiệu điều khiển},\qquad
c(t)\text{ hoặc }y(t):\text{ tín hiệu ra}.
\]

% ======================================================================
\subsection*{CHƯƠNG 1 -- TỔNG QUAN VỀ HỆ THỐNG ĐIỀU KHIỂN TỰ ĐỘNG}
% ======================================================================

\subsubsection*{1. Khái niệm cốt lõi}

Điều khiển là quá trình \textbf{thu thập thông tin, xử lý thông tin và tác động lên hệ thống}
để đáp ứng thực tế gần với mục đích định trước. Điều khiển tự động là quá trình điều khiển
không cần con người trực tiếp can thiệp trong quá trình vận hành.

\begin{center}
\begin{tabular}{|p{0.20\linewidth}|p{0.27\linewidth}|p{0.43\linewidth}|}
\hline
\textbf{Ký hiệu} & \textbf{Tên} & \textbf{Ý nghĩa} \\ \hline
$r(t)$ & Tín hiệu chuẩn / tín hiệu đặt & Giá trị mong muốn của hệ thống \\ \hline
$c(t)$, $y(t)$ & Tín hiệu ra & Giá trị thực tế của đại lượng cần điều khiển \\ \hline
$c_{ht}(t)$ & Tín hiệu hồi tiếp & Tín hiệu đo được đưa trở lại bộ so sánh \\ \hline
$e(t)$ & Sai lệch & Thường có dạng $e(t)=r(t)-c_{ht}(t)$ với hồi tiếp âm \\ \hline
$u(t)$ & Tín hiệu điều khiển & Tín hiệu từ bộ điều khiển tác động lên đối tượng \\ \hline
\end{tabular}
\end{center}

\subsubsection*{2. Hệ vòng hở và hệ vòng kín}

\textbf{Hệ vòng hở (open-loop):} không dùng tín hiệu ra để hiệu chỉnh tác động điều khiển:
\[
r(t)\longrightarrow \text{Bộ điều khiển}\longrightarrow \text{Đối tượng}\longrightarrow c(t).
\]

\textbf{Hệ vòng kín (closed-loop):} dùng tín hiệu phản hồi để so sánh với tín hiệu đặt,
tạo sai lệch và hiệu chỉnh liên tục:
\[
e(t)=r(t)-c_{ht}(t).
\]

Gợi ý nhận biết: nếu sơ đồ có đường từ đầu ra quay trở lại nút so sánh thì đó là
hệ vòng kín; nếu không có đường phản hồi thì là hệ vòng hở.

\subsubsection*{3. Các bài toán điều khiển cơ bản}

\begin{itemize}
  \item \textbf{Phân tích hệ thống:} đã biết cấu trúc và thông số; cần tìm đáp ứng, đặc tính và đánh giá chất lượng.
  \item \textbf{Thiết kế hệ thống:} đã biết đối tượng; cần thiết kế bộ điều khiển để đạt yêu cầu.
  \item \textbf{Nhận dạng hệ thống:} chưa biết đầy đủ cấu trúc/thông số; dùng dữ liệu vào--ra để xác định mô hình.
\end{itemize}

\subsubsection*{4. Các nguyên tắc và phân loại thường gặp}

Trong tài liệu, các nguyên tắc được nhấn mạnh gồm thông tin phản hồi, đa dạng tương xứng,
bổ sung ngoài, dự trữ, phân cấp và cân bằng nội. Ba cấu trúc điều khiển theo thông tin phản hồi
thường gặp là bù nhiễu, san bằng sai lệch và phối hợp hai nguyên tắc trên.

Các cách phân loại chính:
\begin{itemize}
  \item Tuyến tính / phi tuyến.
  \item Bất biến theo thời gian / biến đổi theo thời gian.
  \item Liên tục / rời rạc / số.
  \item SISO (một vào--một ra) / MIMO (nhiều vào--nhiều ra).
  \item Theo mục tiêu: ổn định hoá, theo chương trình, bám theo tín hiệu, thích nghi, tối ưu.
\end{itemize}

\subsubsection*{5. Quy trình thiết lập một hệ thống điều khiển}

Có thể ghi nhớ theo chuỗi:
\[
\boxed{
\text{Yêu cầu kỹ thuật}
\rightarrow
\text{Mô tả hệ vật lý}
\rightarrow
\text{Sơ đồ chức năng}
\rightarrow
\text{Mô hình toán}
\rightarrow
\text{Sơ đồ khối/graph/KGTT}
\rightarrow
\text{Rút gọn}
\rightarrow
\text{Phân tích, thiết kế}
}
\]

\subsubsection*{6. Dạng bài Chương 1: nhận dạng thành phần của hệ}

\textbf{Quy trình:}
\begin{enumerate}
  \item Xác định đại lượng mong muốn $\Rightarrow r(t)$.
  \item Xác định đại lượng thực tế cần điều khiển $\Rightarrow c(t)$.
  \item Xác định thiết bị đo tạo tín hiệu phản hồi.
  \item Xác định bộ điều khiển.
  \item Xác định cơ cấu chấp hành.
  \item Xác định đối tượng điều khiển.
  \item Kiểm tra có đường phản hồi hay không để kết luận vòng hở/vòng kín.
\end{enumerate}

\textbf{Ví dụ mẫu -- điều khiển mức chất lỏng.}
Mức mong muốn là $r(t)$; cảm biến/phao đo mức tạo phản hồi; bộ điều khiển tính sai lệch;
van là cơ cấu chấp hành; bồn là đối tượng; mức thực tế là $c(t)$. Nếu mức đo được quay về
bộ so sánh thì hệ là vòng kín.

\medskip
\noindent
\fbox{\parbox{0.96\linewidth}{
\textbf{Lỗi hay gặp Chương 1:}
nhầm tín hiệu điều khiển $u(t)$ với tín hiệu đặt $r(t)$;
nhầm cảm biến với bộ điều khiển; thấy có cảm biến nhưng không kiểm tra xem tín hiệu đo
có thực sự quay về bộ so sánh hay không.
}}

